\section{Limits of the method}\label{sec:limits}

The decompositions above are produced by rules: contract a shortest cycle (Theorems~N3 and~U3), a shortest directed cycle
of in-degree-two vertices (Theorem~T), or a whole block (Theorem~PC). This section shows where such rules stop.
\Cref{prop:no-private} and \cref{ex:D9} are formalized, by kernel evaluation. The remarks report computations,
labelled as numerical, a cited formula and a heuristic; where parts of them are Lean-verified, they say so, and none of
them is used in a proof. The computations are reproduced by the script \leanfile{tools/verify_remarks.py}, described in
\leanfile{tools/README.md}.

\subsection{Private pieces need not exist}\label{sec:no-private}

The induction of Theorem~N3 removes a piece whose outside parents are private. In directed relations of degree three
such pieces need not exist, even when the half bound holds.

\begin{proposition}[No private proper pieces]\label{prop:no-private}
\formal{Limits.lean}{cay9, no_private_piece}
Let $N$ be the relation of $\Cay(\Z_9,\{1,-1,3\})$, so $N(q)=\{q-1,q+1,q-3\}$ and the children of $q$ are $q+1$, $q-1$
and $q+3$. There is no set $H$ with $\emptyset\ne H\ne\Z_9$ such that
\begin{enumerate}[label=(\alph*)]
\item every vertex of $H$ has at most one parent outside $H$, and
\item every vertex outside $H$ is a parent of at most one vertex of $H$.
\end{enumerate}
\end{proposition}

\begin{outline}
The formal proof evaluates the statement for all $512$ subsets $H$ of $\Z_9$ in the kernel (\texttt{decide +kernel}).
The following argument explains the result. Suppose $H$ exists. Since $\emptyset\ne H\ne\Z_9$, some $a\in H$ has
$a-1\notin H$. By~(a) at $a$, two of its parents $a-1,a+1,a-3$ lie in $H$, so $a+1\in H$ and $a-3\in H$. The children of
$a-1\notin H$ are $a-2$, $a$ and $a+2$; since $a\in H$, (b) gives $a-2\notin H$ and $a+2\notin H$. Now the parents of
$a+1$ are $a$, $a+2$ and $a-2$, and only $a$ lies in $H$, contradicting~(a). The residues $a-3,\dots,a+2$ are distinct
modulo~$9$, so the argument is valid.
\end{outline}

Every piece of a decomposition in context $\emptyset$ satisfies~(a), and a piece whose outside parents are private
satisfies~(b). So on this relation a decomposition must start with the whole vertex set or with a piece whose outside
parents are shared. Both happen. Corollary~CC gives a single piece; explicitly, the records $\{0,4\}$, $\{1,3,6,8\}$,
$\{2,5,7\}$ with the list $0,1,2,3,4,7,5,6,8$ form a global elimination certificate (\lean{z9_elimOn}), which gives
\eqref{eq:star} directly (\lean{rateBound_of_elimOn_univ}, \lean{rateBound_cay9}). And nested charges give a proper
first piece, as follows.

\begin{example}[A proper piece with shared outside parents]\label{ex:D9}
\formal{Limits.lean}{exH1, ex72Piece, ex72Piece4, ex72Piece78, ex72_decomp}
In the relation of \cref{prop:no-private}, let $H_1=\{0,1,2,3,5,6\}$ with records $X=\{0,2,6\}$, $Y=\{1,3,5\}$ and an
empty record $Z$, and the list $2,3,5,6,0,1$ with supplying records $Y,X,X,Z,Y,Z$. The outside parents are $8$ for the
vertices $2$ and $0$, $4$ for $3$ and $5$, and $7$ for $6$ and $1$; each outside vertex is shared by two steps with the
same supplying record, so the certificate is co-supplied (\cref{lem:cosupply}). For instance, at the step of $5$ the
revealed set is $\{2,3\}$ and $N(5)=\{4,6,2\}$ lies in $\{2,3\}\cup X\cup\{4\}$. The relation that remains,
on $\{4,7,8\}$ with $H_1$ as context, is decomposed by the pieces $\{4\}$ and $\{7,8\}$, with records $\{4\},\emptyset$
and $\{7\},\{8\}$. So $H_1,\{4\},\{7,8\}$ is a decomposition of $N$.
\end{example}

\begin{outline}
The three piece certificates are written out (\lean{ex72Piece}, \lean{ex72Piece4}, \lean{ex72Piece78}); their
conditions, including co-supply (\lean{relCharged_of_cosupply}), are checked by kernel evaluation (\texttt{decide}), and
\lean{ex72_decomp} assembles them.
\end{outline}

\paragraph{Why it matters.} For an impossibility proof, \cref{prop:no-private} shows that the private-piece induction of
Theorem~N3 cannot prove half bounds on directed side information of degree three, even on a single circulant: an argument
must use whole blocks, as in Theorem~PC, or outside parents shared by several steps, which is what the nested charging of
Lemma~E provides.

\subsection{Minimum nonempty cores are not enough}\label{sec:cores}

Call $H\subseteq V$ a \emph{core} if every vertex of $H$ has at most one parent outside $H$. A piece in context
$\emptyset$ is a core, and in a cubic graph the minimum nonempty cores are the shortest cycles, which is why Theorem~N3
can always contract a minimum nonempty core. In directed relations this fails.

\begin{remark}[Numerical: a $15$-vertex relation whose minimum nonempty core is not a piece]\label{rem:X15}
Let $\Ex$ be the relation (\lean{x15}) on $\{0,\dots,14\}$ with
\begin{center}
\small
\begin{tabular}{@{}l*{8}{c}@{}}
\toprule
$q$ & 0 & 1 & 2 & 3 & 4 & 5 & 6 & 7\\
$N(q)$ & $\{5,1,6\}$ & $\{0,2,7\}$ & $\{1,3,8\}$ & $\{2,4,6\}$ & $\{3,5,7\}$ & $\{4,0,8\}$ & $\{14,0,10\}$ & $\{13,4,9\}$\\
\midrule
$q$ & 8 & 9 & 10 & 11 & 12 & 13 & 14 & \\
$N(q)$ & $\{9,14,5\}$ & $\{6,11,7\}$ & $\{9,12\}$ & $\{2,13,12\}$ & $\{14,11\}$ & $\{1,8,11\}$ & $\{3,13,12\}$ & \\
\bottomrule
\end{tabular}
\end{center}
The relation $\Ex$ is loopless and TT-free, and every vertex has at most three parents and at most three children
(\lean{x15_props}). Every nonempty core has at least six vertices, and the only core with six vertices is
$C=\{0,\dots,5\}$ (\lean{x15_cores}). The set $C$ is a core, its cyclically consecutive vertices form digons, and its
outside parents $6$, $7$, $8$ each have two antipodal children: $6\to0,3$, $7\to1,4$ and $8\to2,5$ (\lean{x15_hexagon}).
So $C$ is the unique minimum nonempty core, and, as every vertex has at most three parents, it is a bidirected
$6$-cycle. A computation (\leanfile{tools/verify_remarks.py}) shows that $\Ex$ is strongly connected and not
$2$-dicolourable, and an exhaustive search over the partitions of $C$ into at most seven records, all elimination orders
and all supplying records finds no piece certificate for $C$ in context $\emptyset$, even when the steps charged to an
outside parent are only required to be nested (co-supplied outside parents are a special case). More records cannot
help: $C$ has six vertices, so a certificate with more than seven records has at least two empty records, which can be
merged. So $C$ is not a piece in context $\emptyset$, with any number of records. Nevertheless $\Ex$ has a
decomposition, with a single piece: a global elimination certificate with three records, obtained by giving the vertices
$0,\dots,14$ the records $0,1,0,1,0,1,1,1,0,0,0,1,0,0,2$ and using the list $0,4,5,7,1,9,8,2,11,12,10,6,3,13,14$
(\lean{x15_elimOn}). So \eqref{eq:star} holds for $\Ex$ (\lean{rateBound_x15}), but no rule that starts by contracting a
minimum nonempty core can prove it. The claims that $\Ex$ is strongly connected and not $2$-dicolourable and that $C$ is
not a piece are numerical; the other claims of this remark are Lean-verified as cited.
\end{remark}

In the undirected setting the antipodal pattern cannot occur on a shortest cycle: an outside vertex adjacent to $c_0$ and
$c_3$ on a $6$-cycle closes the $5$-cycle $w,c_0,c_1,c_2,c_3$. In $\Ex$ the arc $6\to3$ is one-way
($0\leftrightarrow6$ is a digon); the $5$-cycle $6,0,1,2,3$ of the underlying graph is not a core, since vertex~$6$ has
the two parents $10$ and $14$ outside it, and no nonempty core has fewer than six vertices (\lean{x15_cores}).
The relation lies outside Theorems~U3 and~T: vertex $6$ has five neighbours in the underlying graph, and the vertices with
three parents include the digons of $C$. Its digon components violate~(H1) of \cref{cor:PC-digon}, since vertex $6$ has
the two one-way parents $10$ and $14$ (\lean{x15_oneWay_parents}); Theorem~PC applies with the whole vertex set as one
block, through the global certificate.

\subsection{Sweeps and global certificates}\label{sec:sweeps}

For a chorded digon cycle $K=\{k_0,\dots,k_{L-1}\}$ (indices modulo $L$) such that $N[K]$ is loopless and TT-free, the
\emph{sweep} from $k_z$ in direction $d=\pm1$ is the elimination list $k_z,k_{z+d},\dots,k_{z+(L-1)d}$ with the finest
records that make the later parents of each vertex lie in one record: the classes of the equivalence relation that
merges these later parents, for every vertex. The sweep is \emph{valid} if no vertex lies in the class of its later
parents; since $N[K]$ is loopless, it is then an elimination certificate. The proof of \cref{prop:CC}\ref{it:mirror}
shows that such a cycle (with $N[K]$ loopless and TT-free) whose chords all have one offset $k$ with
$2k\not\equiv\pm1\pmod L$ has a valid sweep, along the cycle or, for the mirror sweep, in the opposite direction; this includes every loopless TT-free
circulant $\Cay(\Z_L,\{1,-1,k\})$.

\begin{remark}[Numerical: chorded cycles without a valid sweep]\label{rem:sweeps}
An adversarial local search over chord sets found twelve TT-free chorded digon cycles, of lengths $13$, $15$ and $16$,
none of whose $2L$ sweeps is valid (\leanfile{tools/verify_remarks.py}). One of them has $L=13$ and the chords
\[
 0\to8,\ 1\to6,\ 3\to11,\ 4\to0,\ 5\to10,\ 6\to9,\ 7\to3,\ 8\to4,\ 9\to12,\ 10\to1,\ 11\to7,\ 12\to2.
\]
All twelve have elimination certificates with three records. For the cycle displayed, give the vertices
$0,\dots,12$ the records $0,2,0,1,0,2,1,2,1,2,1,2,1$; the list $11,9,7,6,5,4,10,1,2,0,8,3,12$ is then an elimination
certificate, which can be checked by hand. So these cycles are of kind~(c) in Theorem~PC; they only show that sweeps, the
method of \cref{prop:CC}, do not cover all chorded cycles.

All twelve cycles are Lean-verified to be loopless TT-free chorded digon cycles: on $\Z_L$, every vertex $v$ has the
parents $v\pm1$, at most three parents and at most three children (\lean{IsChordedCycle}), so beyond the digons of the
cycle there is at most one chord into and at most one chord out of every vertex. Their global elimination
certificates, which the Lean proofs construct with three records, and hence \eqref{eq:star}, are Lean-verified as well.
The displayed cycle is \lean{c13}: \lean{c13_spec} shows that the parents of each vertex are exactly its two cycle
neighbours and the tails of the displayed chords into it, and \lean{c13_isChordedCycle}, \lean{c13_elimOn} and
\lean{rateBound_c13} give the rest. The other eleven are \lean{cyc1}, \dots, \texttt{cyc11}, each with its
certificate \texttt{cyc$K$\_elimOn}; \lean{twelve_cycles} collects the twelve statements \lean{IsChordedCycle} and
\eqref{eq:star}. Apart from \lean{cyc1}, the declarations of these eleven cycles are not listed in the index
individually; their axioms are checked through \lean{twelve_cycles}. Only the claim that none of the $2L$ sweeps is
valid is numerical.
\end{remark}

\begin{remark}[Numerical: the Petersen graph has no global certificate]\label{rem:petersen}
Let $N$ be the neighbourhood relation of the Petersen graph. In a global elimination certificate of $N$
(\cref{def:elimcert}), written as an elimination list, the later neighbours of each vertex lie in one record, different
from the record of the vertex; when the number of records is not bounded, an elimination order admits a certificate
exactly when these forced equalities and inequalities of records are consistent, since the classes of the forced
equalities can then serve as the records. An exhaustive search over all elimination orders of $N$, which keeps the
forced equalities in a union--find structure and abandons an order as soon as a forced inequality fails, finds no
consistent order (\leanfile{tools/verify_remarks.py}). So $N$ has no global elimination certificate, with any number of
records. The Petersen graph, \lean{MotionDepthCodes.petersenGraph} of the fourth paper's library, is triangle-free
(\lean{petersen_cliqueFree}) and cubic, so Theorem~N3, by contrast, gives it a decomposition (\lean{petersen_decomp}).
Following the proof of Theorem~N3, the first piece is a $5$-cycle, whose outside parents (the spokes) are private, and the
second is the remaining $5$-cycle; this description of the pieces is not part of the formal statement. So
decompositions are strictly stronger than global certificates, and Theorem~D strictly stronger than
\cref{prop:global}.
\end{remark}

\subsection{Odd girth and degree}\label{sec:foldcubes}

\begin{remark}[Cited: folded cubes]\label{rem:foldcube}
The folded $(2k+1)$-cube $FQ_{2k+1}=\Cay(\F_2^{2k},\{e_1,\dots,e_{2k},\mathbf 1\})$, where $\mathbf1$ is the all-ones
vector, has $N=2^{2k}$ vertices, degree $2k+1$ and odd girth $2k+1$; for $k=2$ it is the Clebsch graph. Over $\F_2$, the
matrix $I+A$ has rank $\frac12(2^{2k}-2^k)$: this is the case of even dimension $2k$ of a proposition of Barg,
Schwartz and Yohananov~\cite[Proposition~13]{BSY2024}, whose Cayley graph also has the zero vector as a generator and
therefore has adjacency matrix $I+A$. Its kernel is a binary linear storage code of dimension
$N/2+2^{k-1}$; for $k=2,3$ this is also checked numerically (\leanfile{tools/verify_remarks.py}). For $k=2$ it is
Lean-verified in the fourth paper~\cite[Corollary~6.11]{FourthPaper2026}: its Clebsch graph is the folded $5$-cube
(\lean{MotionDepthCodes.clebschGraph_adj_iff}), and $\ker(I+A)$ has dimension $10=N/2+2$
(\lean{MotionDepthCodes.finrank_ker_clebsch}). So \eqref{eq:star} fails on the folded cubes, whose odd girth equals their
degree, for every $k\ge2$; it holds on other graphs whose odd girth equals their degree, such as the $5$-regular
Andrásfai graph $\Andr(5)=K_{14/5}$, of odd girth five (\cref{thm:circular}). In the symmetric case $\dist_E(L_p,R_p)$ is
the length of the shortest odd closed walk through $p$, so here $b_E(r)=0$ for $r<2k+1$. Hence, for every odd $r\ge3$, an exact half bound that assumes only $b_E(r)=0$ can hold at most for maximum degree
$r+1$: the folded $(r+2)$-cube violates it. For $r=3$ it fails already at maximum degree four (\cref{ex:deg4}), so the
degree bound of Theorem~N3 cannot be raised. The excess $2^{k-1}=\sqrt N/2$ is $o(N)$, which is consistent with the
asymptotic form of the paired-distance question~\cite[Question~9.1]{DirectedRecovery2026}.
\end{remark}

\begin{remark}[Heuristic: degree four]\label{rem:degree4}
In the undirected setting, the method needs each step to have at most one outside parent. If a vertex $v$ of a shortest
cycle has two outside neighbours $p$ and $p'$, the step that eliminates it costs $I(v;\{p,p'\}\mid Q)$, which is not a sum
of charges on $p$ and $p'$ that telescope separately. This limitation of Lemma~E is not an artefact: the half bound itself
fails at degree four, on a triangle-free graph of odd girth five (\cref{ex:deg4}), while \cref{thm:circular} covers graphs
of circular chromatic number below three at every degree. The
fourth paper shows that the local linear-algebraic necessary conditions on a smallest counterexample, checked inside
tree-like regions, cannot give rate $\frac12+o(1)$ at degree four or more~\cite[Proposition~6.28]{FourthPaper2026}.
\end{remark}

\begin{example}[The half bound fails at degree four]\label{ex:deg4}
\formal{DegreeFour.lean}{g13, g13_triangleFree, g13_five_cycle, g13_card_nbhd, g13code_isStorageCode, g13code_card,
g13_not_rateBound}
Let $G_{13}$ be the graph on $\{0,\dots,12\}$ in which the neighbours of $0,1,\dots,12$ are, in this order,
$\{4,7,9,10\}$, $\{5,7,8,10\}$, $\{6,7,8,9\}$, $\{8,9,10,11\}$, $\{0,8,11,12\}$, $\{1,9,11,12\}$, $\{2,10,11,12\}$,
$\{0,1,2,12\}$, $\{1,2,3,4\}$, $\{0,2,3,5\}$, $\{0,1,3,6\}$, $\{3,4,5,6\}$ and $\{4,5,6,7\}$. It is $4$-regular and
triangle-free, and $0,4,8,1,7$ is a $5$-cycle, so its odd girth is five. Over the field $\F_4=\{0,1,\omega,\omega^2\}$,
$\omega^2=\omega+1$, let $C\subseteq\F_4^{13}$ be the row space of the $7\times13$ matrix $[I_7\mid D]$, whose
column~$v$ belongs to vertex~$v$, where the rows of~$D$ are
\[
  (1,0,1,1,0,1),\ (1,1,0,1,0,1),\ (1,1,1,0,0,1),\ (0,1,1,1,1,0),\ (1,1,0,0,\omega,\omega),\ (1,0,1,0,\omega,\omega),\
  (1,0,0,1,\omega,\omega).
\]
Every column of $[I_7\mid D]$ is an $\F_4$-linear combination of the columns of the four neighbours of its vertex, so
every symbol of a word of~$C$ is a function of the symbols at its neighbours: $C$ is a storage code for $N_{G_{13}}$ over the
alphabet~$\F_4$, with $\card C=4^7$. Hence $\card C^2=4^{14}>4^{13}=\card\F_4^{\card V}$, and by \cref{thm:codes} the relation
$N_{G_{13}}$ does not satisfy \eqref{eq:star}.
\end{example}

\begin{outline}
The formal proof writes a symbol $a+b\omega$ of $\F_4$ as the bit pair $(a,b)$, so that $C$ is the image of a
$\Z/2$-linear map on $(\Z/2)^{14}$ (\lean{g13enc}, \lean{g13code}); the generator and the recovery coefficients, $2\times2$
blocks for multiplication by elements of $\F_4$, are produced by \leanfile{tools/degree_four.py}. The kernel checks that
each of the fourteen basis words satisfies the recovery identity at every vertex (\lean{g13gen_rec}); linearity extends it
to every word (\lean{g13enc_rec}), which is the storage property (\lean{g13code_isStorageCode}). The first seven symbols of
a word are its message (\lean{g13gen_id}), so the encoding is injective and $\card C=2^{14}$ (\lean{g13code_card}). With
\cref{thm:codes} this contradicts the half-rate bound (\lean{g13_not_rateBound}). Triangle-freeness, the $5$-cycle and the
degrees are decided by the kernel (\lean{g13_triangleFree}, \lean{g13_five_cycle}, \lean{g13_card_nbhd}). The alphabet
of the formal statement is $(\Z/2)^2$, of size four, as for~$\F_4$.
\end{outline}
